% =====================================================================
% pertemuan-15.tex
% Pertemuan 15 - AST, Validasi Kode, dan Penanganan Kesalahan
% Sub-CPMK 4.4, 1.7 - Praktik terbimbing, penutup proyek
% =====================================================================
\documentclass[aspectratio=169,10pt,t]{beamer}
\usepackage{config/beamerthemeubb}

\UBBmeeting{15}

\begin{document}

\UBBtitleframe{15}{AST, Validasi Kode, dan Penanganan Kesalahan}{Praktik terbimbing \textperiodcentered{} Penutup proyek}

% ---------------------------------------------------------------------
\begin{frame}{Capaian Pertemuan}
  \textbf{CPMK induk:} CPMK 04 --- Mengimplementasikan parser sederhana
  untuk validasi kode; CPMK 01 --- Menjelaskan tahapan kompilasi.

  \vspace{2pt}
  \begin{block}{Sub-CPMK yang dinilai}
    \footnotesize
    \begin{itemize}
      \setlength{\itemsep}{1pt}
      \item \textbf{Sub-CPMK 4.4:} Membuat AST dan validator kode
            untuk subset bahasa C, lalu menguji keluaran pesan
            kesalahan sintaksis yang informatif.
      \item \textbf{Sub-CPMK 1.7:} Menguji dan men-\textit{debug}
            penanganan kesalahan semantik sehingga program menghasilkan
            pesan kesalahan yang menyebut baris, nama, dan jenis
            pelanggaran.
    \end{itemize}
  \end{block}
\end{frame}

% ---------------------------------------------------------------------
\begin{frame}{Agenda Pertemuan (150 menit)}
  \small
  \begin{tabular}{@{}p{1.4cm}p{12.3cm}@{}}
    \toprule
    \textbf{Menit} & \textbf{Kegiatan} \\
    \midrule
    20 & Pembuka: posisi M5 dan cakupan proyek akhir \\
    25 & Ceramah: desain AST dan traversal \\
    30 & Praktik: validator dan aturan validasi \\
    30 & Praktik: diagnostik dan pengujian pesan \\
    25 & Pengujian akhir pipeline dan demo \\
    20 & Laporan, refleksi, dan penutup proyek \\
    \bottomrule
  \end{tabular}

  \vspace{6pt}
  {\scriptsize\color{ubbblue!75}Porsi praktik 85 menit. Pertemuan ini
  menutup increment proyek akhir (M5) sebelum UAS.}
\end{frame}

% ---------------------------------------------------------------------
\begin{frame}{Posisi Pertemuan pada Rantai Proyek}
  \UBBchain{6}

  \vspace{2pt}
  \begin{itemize}
    \item Pertemuan ini menutup simpul \textbf{AST} dan menyatukan
          seluruh rantai menjadi validator.
    \item Pipeline lengkap: lexer, parser, type checker, validator,
          diagnostik.
    \item Keluaran akhir adalah CLI \texttt{scc} untuk Subset-C.
  \end{itemize}
\end{frame}

% ---------------------------------------------------------------------
\begin{frame}{Prinsip Desain AST}
  \begin{enumerate}
    \setlength{\itemsep}{2pt}
    \item Setiap simpul merepresentasikan satu konstruksi bahasa.
    \item Informasi sintaksis yang tidak bermakna dihilangkan.
    \item Simpul menyimpan lokasi untuk diagnostik.
    \item Struktur mendukung traversal yang mudah.
    \item Tipe dan informasi semantik dapat ditambahkan tanpa mengubah
          bentuk dasar.
  \end{enumerate}

  \vspace{4pt}
  \begin{block}{Generik versus spesifik}
    \footnotesize
    Simpul generik mudah dipelajari, tetapi kurang \eng{type-safe}.
    Simpul spesifik memberi kejelasan dan mencegah akses field yang
    salah; untuk proyek besar biasanya lebih baik.
  \end{block}
\end{frame}

% ---------------------------------------------------------------------
\begin{frame}[fragile]{Struktur Simpul AST (\file{ast.h})}
  \begin{columns}[T,onlytextwidth]
    \begin{column}{0.58\textwidth}
      \lstinputlisting[language=C,firstline=37,lastline=51,
        firstnumber=37,
        caption={Struktur \texttt{AstNode} dan antarmukanya},
        label={lst:p15-ast}]{../buku_ajar/code/subsetc/include/ast.h}
    \end{column}
    \begin{column}{0.38\textwidth}
      \begin{itemize}
        \item \texttt{kind}: jenis simpul.
        \item \texttt{token}: lokasi dan lexeme.
        \item \texttt{decl\_type}: tipe deklarasi.
        \item \texttt{left}/\texttt{right}: operand.
        \item \texttt{children}/\texttt{next}: daftar anak.
      \end{itemize}
    \end{column}
  \end{columns}
\end{frame}

% ---------------------------------------------------------------------
\begin{frame}{Konstruksi AST}
  \begin{itemize}
    \item Aturan yang hanya mengatur sintaksis tidak membuat simpul.
    \item Aturan operator membuat simpul biner dengan dua anak.
    \item Aturan daftar membuat simpul dan menambahkan anak.
    \item Deklarasi menyimpan tipe pada \texttt{decl\_type}.
    \item AST dibangun saat parsing, dari bawah ke atas.
  \end{itemize}

  \vspace{4pt}
  \begin{exampleblock}{Dump AST untuk \texttt{int a = 2 + 3;}}
    \footnotesize
    \texttt{PROGRAM > FUNC main > BLOCK > DECL a > BINARY + >
    INT\_LIT 2, INT\_LIT 3}.\\
    Tidak ada titik koma atau kurung pada AST.
  \end{exampleblock}
\end{frame}

% ---------------------------------------------------------------------
\begin{frame}{Konstruksi Simpul: \texttt{mknode} dan \texttt{mkleaf}}
  \begin{itemize}
    \item Juarna memakai dua operasi dasar: \texttt{mkleaf} untuk daun
          dan \texttt{mknode} untuk simpul operator.
    \item \texttt{mkleaf(id, a)}, \texttt{mkleaf(id, b)},
          \texttt{mkleaf(id, c)}.
    \item \texttt{mknode(*, b, c)} menghasilkan \(n_1\).
    \item \texttt{mknode(+, a, n\_1)} menghasilkan akar \(n_2\).
    \item Di C, keduanya dibungkus menjadi \texttt{ast\_new}.
  \end{itemize}

  \vspace{4pt}
  \begin{exampleblock}{Hasil untuk \texttt{a + b * c}}
    \footnotesize
    Akar \texttt{+}; anak kiri \texttt{a}; anak kanan \texttt{*};
    cucu \texttt{b} dan \texttt{c}.
  \end{exampleblock}
\end{frame}

% ---------------------------------------------------------------------
\begin{frame}{Traversal AST}
  \begin{columns}[T,onlytextwidth]
    \begin{column}{0.48\textwidth}
      \begin{itemize}
        \item \textbf{Pre-order}: simpul lalu anak-anak.
        \item Cocok untuk analisis yang butuh konteks orang tua,
              misalnya validasi loop.
        \item \textbf{Post-order}: anak-anak lalu simpul.
        \item Cocok untuk sintesis, misalnya perhitungan tipe.
      \end{itemize}
    \end{column}
    \begin{column}{0.46\textwidth}
      \centering
      \begin{tikzpicture}[
          nd/.style={draw=ubbblue!40,fill=ubbblue!7,rounded corners=2pt,
            minimum width=0.7cm,minimum height=0.48cm,font=\scriptsize},
          ar/.style={-{Stealth[length=1.3mm]},draw=ubbblue!55}
        ]
        \node[nd] (r) at (0,0) {\(+\)};
        \node[nd] (a) at (-1.2,-0.8) {\code{a}};
        \node[nd] (m) at (1.2,-0.8) {\(*\)};
        \node[nd] (b) at (0.5,-1.6) {\code{b}};
        \node[nd] (c) at (1.9,-1.6) {\code{c}};
        \draw[ar] (r)--(a); \draw[ar] (r)--(m);
        \draw[ar] (m)--(b); \draw[ar] (m)--(c);
      \end{tikzpicture}
      \par {\scriptsize Pre-order: \(+,a,*,b,c\)}
    \end{column}
  \end{columns}
\end{frame}

% ---------------------------------------------------------------------
\begin{frame}{Pola \textit{Visitor} dan Konteks}
  \begin{itemize}
    \item \eng{Visitor} memisahkan operasi dari struktur data.
    \item Implementasi C: fungsi yang menerima simpul dan konteks.
    \item Proyek Subset-C memakai \texttt{switch} pada \texttt{kind}.
    \item Setiap operasi baru menjadi fungsi traversal terpisah:
          \texttt{typecheck}, \texttt{validate}.
    \item \texttt{visit\_children} menghindari duplikasi kode
          traversal.
  \end{itemize}

  \vspace{4pt}
  \begin{block}{Konteks traversal}
    \footnotesize
    Berisi \texttt{loop\_depth}, \texttt{return\_type}, dan
    \texttt{has\_return}. Konteks diperbarui saat masuk atau keluar
    simpul.
  \end{block}
\end{frame}

% ---------------------------------------------------------------------
\begin{frame}{Aturan Validasi Kode}
  \begin{enumerate}
    \setlength{\itemsep}{2pt}
    \item \texttt{break} dan \texttt{continue} hanya di dalam loop.
    \item Fungsi non-\texttt{void} harus mengembalikan nilai.
    \item \texttt{main} harus \texttt{int main()} tanpa parameter.
    \item \texttt{return} hanya di dalam fungsi.
    \item Pernyataan setelah \texttt{return} tidak terjangkau.
    \item Pemanggilan fungsi sesuai \eng{signature}.
  \end{enumerate}

  \vspace{3pt}
  {\footnotesize\color{ubbblue!80}Setiap aturan memiliki pesan yang
  jelas dan lokasi yang tepat.}
\end{frame}

% ---------------------------------------------------------------------
\begin{frame}{Validasi Multi-Pass}
  \begin{center}
    \begin{tikzpicture}[
        nd/.style={draw=ubbblue!40,fill=ubbblue!7,rounded corners=3pt,
          minimum width=2.6cm,minimum height=0.62cm,font=\scriptsize,
          text=ubbblue},
        data/.style={nd,draw=ubbgold,fill=ubbgold!20},
        ar/.style={-{Stealth[length=1.6mm]},draw=ubbblue!55,thick}
      ]
      \node[nd] (parse) {Parser};
      \node[nd,right=0.45cm of parse] (sem) {Type checker};
      \node[nd,right=0.45cm of sem] (val) {Validator};
      \node[data,below=0.6cm of val] (rep) {Diagnostik};
      \draw[ar] (parse) -- (sem);
      \draw[ar] (sem) -- (val);
      \draw[ar] (val) -- (rep);
    \end{tikzpicture}
  \end{center}

  \begin{enumerate}
    \setlength{\itemsep}{1pt}
    \item Parsing menghasilkan AST.
    \item Analisis semantik membangun tabel simbol dan memeriksa tipe.
    \item Validasi memeriksa aturan kontrol dan kelengkapan.
    \item Pelaporan mengumpulkan dan mengurutkan diagnostik.
  \end{enumerate}
\end{frame}

% ---------------------------------------------------------------------
\begin{frame}[fragile]{Validator: Fungsi \texttt{visit}}
  \begin{columns}[T,onlytextwidth]
    \begin{column}{0.60\textwidth}
      \lstinputlisting[language=C,firstline=31,lastline=45,
        firstnumber=31,
        caption={Awal fungsi \texttt{visit} pada \file{validate.c}},
        label={lst:p15-visit}]{../buku_ajar/code/subsetc/src/validate.c}
    \end{column}
    \begin{column}{0.36\textwidth}
      \begin{itemize}
        \item Simpul fungsi memulai konteks baru.
        \item \texttt{return\_type} dan \texttt{loop\_depth} diatur.
        \item Anak dikunjungi dengan konteks tersebut.
        \item Fungsi non-void tanpa \texttt{return} dilaporkan.
      \end{itemize}
    \end{column}
  \end{columns}
\end{frame}

% ---------------------------------------------------------------------
\begin{frame}[fragile]{Titik Masuk Validator}
  \begin{columns}[T,onlytextwidth]
    \begin{column}{0.56\textwidth}
      \lstinputlisting[language=C,firstline=92,lastline=104,
        firstnumber=92,
        caption={Fungsi \texttt{validate\_program}},
        label={lst:p15-validate}]{../buku_ajar/code/subsetc/src/validate.c}
    \end{column}
    \begin{column}{0.40\textwidth}
      \begin{itemize}
        \item Konteks awal: kedalaman loop 0.
        \item Tipe kembalian awal \texttt{void}.
        \item \texttt{visit\_children} menelusuri program.
        \item Mengembalikan status kesalahan.
        \item Driver menggabungkan hasil semua fase.
      \end{itemize}
    \end{column}
  \end{columns}
\end{frame}

% ---------------------------------------------------------------------
\begin{frame}{Klasifikasi Kesalahan}
  \scriptsize
  \setlength{\tabcolsep}{4pt}
  \begin{tabular}{@{}p{2.6cm}p{3.2cm}p{7.2cm}@{}}
    \toprule
    \textbf{Kelas} & \textbf{Fase} & \textbf{Contoh} \\
    \midrule
    Leksikal & Lexer & karakter ilegal, string tidak ditutup \\
    Sintaksis & Parser & titik koma hilang, kurung tidak seimbang \\
    Semantik & Type checker & identifier tidak dideklarasikan, tipe
      salah \\
    Validasi & Validator & \texttt{break} di luar loop, return hilang \\
    \bottomrule
  \end{tabular}

  \vspace{4pt}
  {\scriptsize\color{ubbblue!75}Kesalahan \textit{run-time} bukan
  kesalahan kompilasi. Guntoro membatasi jenis kesalahan program pada
  leksikal, sintaksis, dan semantik.}
\end{frame}

% ---------------------------------------------------------------------
\begin{frame}{Contoh Kesalahan menurut Guntoro}
  \scriptsize
  \setlength{\tabcolsep}{4pt}
  \begin{tabular}{@{}p{2.3cm}p{5.4cm}p{5.2cm}@{}}
    \toprule
    \textbf{Jenis} & \textbf{Contoh} & \textbf{Penjelasan} \\
    \midrule
    Leksikal & \texttt{whle (a < b)} & kata kunci salah eja \\
    Sintaksis & \texttt{a = x + (b * (c + d);} & kurung tidak seimbang \\
    Semantik & \texttt{int siswa; siswa = "Yanuar";} & tipe string
      ke integer \\
    Semantik & \texttt{b = b + 1;} & variabel belum dideklarasikan \\
    \bottomrule
  \end{tabular}

  \vspace{4pt}
  {\scriptsize\color{ubbblue!75}Pesan menyebut baris, nama yang
  terlibat, dan jenis pelanggaran.}
\end{frame}

% ---------------------------------------------------------------------
\begin{frame}{Pemulihan dan Batas Pelaporan}
  \begin{columns}[T,onlytextwidth]
    \begin{column}{0.48\textwidth}
      \textbf{Strategi pemulihan}
      \begin{itemize}
        \item lexer: lewati karakter ilegal;
        \item parser: sinkronisasi ke \texttt{;} atau \texttt{\}};
        \item type checker: \texttt{TYPE\_ERROR} lalu lanjut;
        \item validator: catat dan lanjut ke simpul berikutnya.
      \end{itemize}
    \end{column}
    \begin{column}{0.48\textwidth}
      \textbf{Batas pelaporan}
      \begin{itemize}
        \item Pemulihan agresif dapat menghasilkan pesan palsu.
        \item Setelah sejumlah kesalahan, berhenti dan beri ringkasan.
        \item Batas dapat dikonfigurasi, misalnya 20 kesalahan.
        \item Fokus pada akar masalah.
      \end{itemize}
    \end{column}
  \end{columns}

  \vspace{4pt}
  {\footnotesize\color{ubbblue!80}Pemulihan memungkinkan beberapa
  kesalahan dilaporkan dalam satu proses.}
\end{frame}

% ---------------------------------------------------------------------
\begin{frame}[fragile]{Format Diagnostik dan Kualitas Pesan}
  \begin{block}{Format konsisten}
    \centering
    \texttt{berkas:baris:kolom: error: pesan}
  \end{block}

  \begin{columns}[T,onlytextwidth]
    \begin{column}{0.48\textwidth}
      \begin{lstlisting}[language=C,numbers=none,
        caption={Contoh keluaran diagnostik}]
tests/invalid/break_outside.sc:2:5:
  error: break hanya valid di dalam loop
      \end{lstlisting}
    \end{column}
    \begin{column}{0.48\textwidth}
      \begin{itemize}
        \item buruk: \texttt{error};
        \item cukup: \texttt{baris 4: type error};
        \item baik: \texttt{baris 4, kolom 9: tidak dapat menetapkan
              float ke int; gunakan cast eksplisit}.
      \end{itemize}
    \end{column}
  \end{columns}
\end{frame}

% ---------------------------------------------------------------------
\begin{frame}[fragile]{Contoh: \texttt{break} di Luar Loop}
  \begin{columns}[T,onlytextwidth]
    \begin{column}{0.48\textwidth}
      \begin{lstlisting}[language=C,numbers=none,
        caption={Program tidak valid}]
int main() {
    break;
    return 0;
}
      \end{lstlisting}
    \end{column}
    \begin{column}{0.48\textwidth}
      \begin{itemize}
        \item Validator memasuki fungsi dengan
              \texttt{loop\_depth = 0}.
        \item Simpul \texttt{BREAK} diperiksa.
        \item Kedalaman nol: laporkan
              \texttt{break hanya valid di dalam loop}.
        \item Baris dan kolom menunjuk token \texttt{break}.
      \end{itemize}
    \end{column}
  \end{columns}
\end{frame}

% ---------------------------------------------------------------------
\begin{frame}[fragile]{Contoh: Fungsi Tanpa \texttt{return}}
  \begin{columns}[T,onlytextwidth]
    \begin{column}{0.48\textwidth}
      \begin{lstlisting}[language=C,numbers=none,
        caption={Fungsi non-void tanpa return}]
int hitung(int x) {
    int y = x + 1;
}
      \end{lstlisting}
    \end{column}
    \begin{column}{0.48\textwidth}
      \begin{itemize}
        \item Tipe fungsi \texttt{int}; \texttt{has\_return = 0}.
        \item Setelah anak dikunjungi, validator memeriksa flag.
        \item Pesan:
              \texttt{fungsi non-void harus mengembalikan nilai}.
        \item Lokasi menunjuk nama fungsi.
      \end{itemize}
    \end{column}
  \end{columns}
\end{frame}

% ---------------------------------------------------------------------
\begin{frame}{Kegiatan Praktikum: M5 Proyek Akhir}
  \begin{block}{Langkah kerja}
    \footnotesize
    \begin{enumerate}
      \setlength{\itemsep}{1pt}
      \item Integrasikan seluruh modul: lexer, parser, type checker,
            validator.
      \item Tambahkan aturan validasi yang belum tercakup.
      \item Uji dengan program valid dan tidak valid.
      \item Ukur waktu dan penggunaan memori.
      \item Tulis laporan dan siapkan demo.
    \end{enumerate}
  \end{block}

  \vspace{3pt}
  \begin{exampleblock}{Keluaran M5}
    \footnotesize
    CLI \texttt{scc}; pipeline lengkap; minimal 40 kasus uji; pesan
    kesalahan yang menyebut baris, nama, dan jenis pelanggaran; log
    debug; laporan dan demo.
  \end{exampleblock}
\end{frame}

% ---------------------------------------------------------------------
\begin{frame}{Tugas dan Artefak}
  \begin{block}{Artefak Sub-CPMK 4.4 dan 1.7}
    AST, validator kode, dan pesan kesalahan yang menyebut baris, nama,
    serta jenis pelanggaran, disertai keluaran pengujian.
  \end{block}

  \small
  \begin{itemize}
    \item CLI memproses program valid tanpa kesalahan.
    \item CLI menolak program tidak valid dengan pesan jelas.
    \item Diagnostik menyebut baris, nama, dan jenis pelanggaran.
    \item Minimal 40 kasus uji valid dan tidak valid.
    \item Log \eng{debug} dan perbaikan terdokumentasi.
    \item Laporan proyek dan demo selesai.
  \end{itemize}

  \vspace{3pt}
  {\footnotesize\color{ubbblue!75}Pertemuan ini menutup increment
  proyek akhir; UAS menilai demonstrasi proyek dan ujian tulis.}
\end{frame}

% ---------------------------------------------------------------------
\begin{frame}{Rangkuman dan Penutup Proyek}
  \begin{itemize}
    \item AST yang baik ringkas, menyimpan lokasi, dan mudah
          ditelusuri.
    \item Traversal pre-order dan post-order melayani kebutuhan
          konteks dan sintesis.
    \item Validasi multi-pass memeriksa aturan kontrol setelah tipe.
    \item Kesalahan diklasifikasikan menurut fase dan dipulihkan
          secara terkendali.
    \item Diagnostik yang baik menyebut berkas, baris, kolom, nama,
          dan jenis pelanggaran.
  \end{itemize}

  \vspace{4pt}
  \begin{exampleblock}{Menuju UAS}
    Proyek akhir didemonstrasikan pada Pertemuan 16. Pastikan pipeline
    berjalan pada berkas uji, log debug lengkap, dan laporan siap.
  \end{exampleblock}
\end{frame}

\end{document}
